% ===== tr-01 块1：第 2 章 pp.1--5（原文页 1--5）=====
% [原文 p.1]
\section{镜像对称：SYZ 之前}
Calabi--Yau 流形在数学中之所以重要，原因有很多：它们给出了度量 Ricci 平坦但非平坦的 Einstein 流形的例子，并作为一类 Kodaira 维数为 0 的代数簇，在代数几何中占据特殊地位。但是，当 Yau [177, 178] 于 1976 年给出 Calabi 猜想的著名解答时，没有人料到这些流形还会在物理学中、更确切地说在弦论 [13] 中扮演如此不可或缺的角色------弦论是统一广义相对论与量子场论的一个候选理论，而它的发展反过来又导致了一个被称为镜像对称的数学现象的发现，该现象引发了大量研究，并深刻影响了许多数学分支。

故事始于 20 世纪 80 年代末，当时 Dixon [44] 与 Lerche--Vafa--Warner [119] 发现，弦论或许并不区分 Calabi--Yau 流形；更确切地说，两个不同的 Calabi--Yau 流形在用作时空隐藏额外维度的紧化时，可能给出等价的弦论。这一惊人的想法很快被 Greene--Plesser [67] 与 Candelas--Lynker--Schimmrigk [14] 各自实现，他们独立找到了这类 Calabi--Yau 流形对的非平凡例子；它们被称为镜像对，因为这样的两个流形表现出 Hodge 数的互换。

数学家真正对镜像对称产生兴趣，是在 Candelas、de la Ossa、Green 与 Parkes 利用它完成一项引人注目的计算之后：他们计算了 $\PP^4$ 中五次三维流形上的有理曲线条数 [12]，解决了枚举几何中一个长期悬而未决的问题，并在当时引起了代数几何学家的广泛关注。

这触发了 Gromov--Witten 理论等重要学科的发展，并最终导致 Givental [63, 64] 与 Lian--Liu--Yau [125, 126, 127, 128] 各自独立证明了镜像定理；这些工作特别表明，Candelas 等人的计算在数学上是正确的。这标志着镜像对称数学研究的第一个里程碑。接下来的问题是：如何以一种内在的、数学的方式来理解镜像对称。

第一个这样的表述是 Kontsevich 的同调镜像对称（homological mirror symmetry，HMS）猜想，它在其 1994 年国际数学家大会（ICM）报告 [111] 中提出。用物理学的术语来说，一个 Calabi--Yau 流形 $X$ 决定了两个拓扑弦理论：A-模型与 B-模型，它们分别由 $X$ 的辛几何与复几何所控制 [169, 175]。从这个观点看，镜像对称可以理解为 $X$ 的 A-模型（辛几何）与其镜像 $\check X$ 的 B-模型（复几何）之间的同构，反之亦然。上面关于五次三维流形的枚举预测，就是这个更大图景的众多有趣体现之一。

% [原文 p.2]
Kontsevich 的 HMS 猜想把镜像对称简洁地表述为：$X$ 中拉格朗日子流形的 Fukaya 范畴（A-模型）与镜像 $\check X$ 上凝聚层的导出范畴（B-模型）之间的等价。他的猜想既深刻又优美，并被认为蕴含镜像对称的枚举预测。然而，它并未指出如何找到这样的等价，也没有告诉我们如何构造给定 Calabi--Yau 流形的镜像。
% [Tutor Note: 原文 p.1 末句 "The above enumerative predictions about the quintic 3-fold is one of many interesting manifestations..." 主语为复数 predictions 而谓语用 is，语法上不合；译文按中文习惯处理，不改变数学内容。]
% [Tutor Note: 镜像流形记号：原文通篇用 $\check X$（X 上加 check），不是 $\bar X$，也不是 $\tilde X$；本译本统一保留 $\check X$。]

\section{SYZ 猜想的表述}
1996 年，Strominger、Yau 与 Zaslow [163] 提出了一个开创性的构想，以几何的方式揭示了一对镜像 Calabi--Yau 流形之间的密切关系：

\begin{conjecture}{2.2.1}[SYZ 猜想，见文献 {[163]}]
设 $X$ 与 $\check X$ 是互为镜像的 Calabi--Yau 流形，则
\begin{enumerate}
\item[(i)] $X$ 与 $\check X$ 都容许带截面的特殊拉格朗日环面纤维化 $\mu:X\to B$ 与 $\check\mu:\check X\to B$，它们以同一个底 $B$ 为底：
\[
\begin{tikzcd}
X \arrow[dr, "\mu"] & & \check X \arrow[dl, "\check\mu"] \\
& B &
\end{tikzcd}
\]
并且这两个纤维化在纤维意义下互为对偶，即：对 $b\in B$，若纤维 $\mu^{-1}(b)\subset X$ 与 $\check\mu^{-1}(b)\subset\check X$ 非奇异，则它们是对偶环面；并且
\item[(ii)] 存在纤维方向上的 Fourier 型变换，它们负责在 $X$ 上的辛几何数据（相应地，复几何数据）与 $\check X$ 上的复几何数据（相应地，辛几何数据）之间进行互换。
\end{enumerate}
\end{conjecture}
% [Tutor Note: 上述图中的两条箭头都指向同一个底 $B$：$\mu$ 与 $\check\mu$ 是同一底上的两个纤维化，"fiberwise dual"（纤维对偶）指逐纤维取对偶环面。]
% [Tutor Note: 原文 (ii) 的 "symplectic-geometric (resp. complex-geometric) data on X and complex-geometric (resp. symplectic-geometric) data on $\check X$" 中两个 resp. 是交叉对应的，译文用"相应地"两次标出，不可读成同序对应。]

简而言之，这表明神秘的镜像现象可以被简单地理解为一种 Fourier 变换，即所谓的 T 对偶。这个非凡而影响深远的猜想不仅为镜像对称提供了优美的几何解释，还给出了一个具体的镜像构造：任给 Calabi--Yau 流形 $X$，其镜像由 $X$ 上某个特殊拉格朗日环面纤维化在纤维方向取对偶而得到。它随即成为镜像对称数学研究中的一条主要途径（另一条是 Kontsevich 的 HMS 猜想），并吸引了数学家与物理学家的广泛关注。

我们简要回顾 SYZ 猜想背后的启发式推理。首先，弦论中的一个关键思想是 Dirichlet 膜（D-膜）的存在。物理论证表明，B-模型中的 D-膜（或简称 B-膜）是复子簇上的凝聚层，而 A-模型中的 D-膜（A-膜）是带平坦联络的特殊拉格朗日子流形。由于镜像对称预言 $X$ 的 A-模型与 $\check X$ 的 B-模型之间的等价，$X$ 上一个 A-膜的模空间应当与 $\check X$ 上镜像 B-膜的模空间等同。

注意 $\check X$ 中的一个点就是一个 B-膜，其模空间就是 $\check X$ 本身，因此它应当与 $X$ 上某个 A-膜 $(L,\nabla)$ 的模空间等同，其中 $L\subset X$ 是特殊拉格朗日子流形，$\nabla$ 是 $L$ 上的平坦 $U(1)$-联络。由于 $\check X$ 被它的点所扫过，$X$ 也应当被 $L$ 的形变所扫过。

% [原文 p.3]
现在，McLean 定理 [136] 告诉我们，特殊拉格朗日子流形 $L\subset X$ 的模空间是无障碍的，并以 $H^1(L;\R)$ 为模型；而 $L$ 上平坦 $U(1)$-联络的模空间（模去规范变换）由 $H^1(L;\R)/H^1(L;\Z)$ 给出。因此为了维数匹配，应当有 $\dim H^1(L;\R)=\dim_{\C}\check X=n$，从而 $X$ 应当容许一个特殊拉格朗日环面纤维化
\[
\mu:X\to B.
\]
% [Tutor Note: 维数匹配的关键：McLean 定理给出 $\dim_{\R}H^1(L;\R)=n$，而平坦 $U(1)$-联络的模空间 $H^1(L;\R)/H^1(L;\Z)$ 是 $n$ 维环面；要让一个点的模空间（0 维）与它等同，需先固定 $L$ 的形变，故纤维 $L$ 必须是 $n$ 维环面。]

此外，流形 $\check X$ 本身可以看作一个 B-膜，其模空间是单点，并且它与 $\check X$ 中每个点恰好相交一次，因此对应的 A-膜 $(L,\nabla)$ 应当给出 $\mu$ 的一个特殊拉格朗日截面 $\sigma$，且 $H^1(\sigma;\R)=0$。特别地，底 $B$ 的第一 Betti 数应当为 $b_1=0$。

对 $\check X$ 应用同样的论证，得到带截面的特殊拉格朗日环面纤维化
\[
\check\mu:\check X\to\check B.
\]

对环面纤维 $L_b:=\mu^{-1}(b)\subset X$，其对偶 $L_b^\vee$ 可以看作 $L$ 上平坦 $U(1)$-联络的模空间；在镜像对称下，这些联络对应于 $\check X$ 中的点。这表明 $L_b^\vee$ 是 $\check X$ 中的子流形。更精细的论证表明，$L_b^\vee$ 应当与 $\check\mu$ 的一个特殊拉格朗日环面纤维等同，于是我们推出 $\mu$ 与 $\check\mu$ 是同一底上的纤维化，且在纤维意义下互为对偶。
% [Tutor Note: 原文此处写 "the moduli space of flat $U(1)$-connections on $L$"，但本句开头讨论的是纤维 $L_b:=\mu^{-1}(b)$；这里的 $L$ 应理解为该纤维 $L_b$（原文混用了两个记号），不含额外数学内容。]
% [Tutor Note: 内乘记号：原文的 $\iota_v$（iota）作用在形式上下标为 $v$；不要误抄成 $\imath$ 或 $i$。]

注意上述论证给出了一个变换，它把 A-膜 $(L,\nabla)$（其中 $L$ 是 $\mu$ 的纤维）送到 $\check X$ 中的点（作为 B-膜）。这是所谓 SYZ 变换的一个实例；SYZ 变换是 Fourier 型变换，把 $X$ 上的辛几何数据映为 $\check X$ 上的复几何数据。

在原始的 SYZ 论文 [163] 中，人们注意到镜像 $\check X$ 上的 Ricci 平坦度量的行为应当不同于 Greene--Shapere--Vafa--Yau [68] 早先构造的半平坦 Calabi--Yau 度量。因此，SYZ 镜像构造一般而言应当被瞬子修正或量子修正所改变。正如我们将看到的，研究镜像对称的一个关键步骤就是理解这些修正；它们应当来自 SYZ 变换的高阶 Fourier 模。

\section{半平坦 SYZ}
当特殊拉格朗日环面纤维化中没有奇异纤维时，SYZ 构造可以被很好地实现。首先，McLean 的经典结果 [136] 给出特殊拉格朗日环面纤维化 $\mu:X\to B$ 的底流形 $B$ 上两个自然定义的整仿射结构\footnote{流形上的整仿射结构是指一个图册，其转移映射都是整仿射线性变换。}：辛仿射结构与复仿射结构。更具体地说，纤维 $L_b:=\mu^{-1}(b)$ 的一个法向量场 $v$ 决定一个 1-形式 $\alpha:=-\iota_v\omega\in\Omega^1(L_b;\R)$ 与一个 $(n-1)$-形式 $\beta:=\iota_v\operatorname{Im}\Omega\in\Omega^{n-1}(L_b;\R)$，其中 $\omega$ 与 $\Omega$ 分别是 $X$ 上的 Kähler 形式与全纯体积形式。McLean [136] 证明了：相应的形变是特殊拉格朗日形变，当且仅当 $\alpha$ 与 $\beta$ 都闭。用 $\alpha$（相应地 $\beta$）的上同调类把 $TB$ 与 $H^1(L_b;\R)$（相应地 $H^{n-1}(L_b;\R)$）等同起来，我们就得到 $B$ 上的辛（相应地复）仿射结构。这还给出 $B$ 上的 McLean 度量
\[
g(v_1,v_2):=-\int_{L_b}\iota_{v_1}\omega\wedge\iota_{v_2}\operatorname{Im}\Omega .
\]
% [Tutor Note: 脚注 1 的"整仿射线性变换"即形如 $x\mapsto Ax+b$ 的变换，其中 $A\in GL(n;\Z)$；这是整仿射结构与光滑仿射结构的区别所在（后者只要求 $A\in GL(n;\R)$）。]
% [Tutor Note: $\beta=\iota_v\operatorname{Im}\Omega$ 是 $(n-1)$-形式，故与 $\Omega^{n-1}(L_b;\R)$ 匹配；若误抄成 $\Omega^{n-2}$ 则维数不匹配。]

% [原文 p.4]
Hitchin 在其优美的论文 [95] 中利用这些结构与 Legendre 变换来阐释 SYZ 猜想。设 $x_1,\dots,x_n$ 是 $B$ 上关于辛仿射结构的局部坐标，则 McLean 度量可以写成 $B$ 上凸函数 $\phi$ 的 Hessian，即 $g\bigl(\frac{\partial}{\partial x_i},\frac{\partial}{\partial x_j}\bigr)=\frac{\partial^2\phi}{\partial x_i\partial x_j}$。令 $\check x_i:=\partial\phi/\partial x_i$（$i=1,\dots,n$），恰好给出 $B$ 上关于复仿射结构的坐标；而若
\[
\check\phi:=\sum_{i=1}^{n}\check x_ix_i-\phi(x_1,\dots,x_n)
\]
是 $\phi$ 的 Legendre 变换，则我们有 $x_i=\partial\check\phi/\partial\check x_i$ 以及 $g\bigl(\frac{\partial}{\partial\check x_i},\frac{\partial}{\partial\check x_j}\bigr)=\frac{\partial^2\check\phi}{\partial\check x_i\partial\check x_j}$。
% [Tutor Note: Legendre 变换的符号：对偶坐标记作 $\check x_i$、对偶势记作 $\check\phi$；与原坐标 $x_i$、原势 $\phi$ 一一对应。$\check\phi$ 的定义式是 $\sum_i\check x_ix_i-\phi$，不是 $\phi-\sum_i\check x_ix_i$。]

若纤维化 $\mu:X\to B$ 容许一个拉格朗日截面，则 Duistermaat 的一个定理 [45] 会给出整体作用-角坐标，使得在辛几何意义下
\[
X=T^*B/\Lambda^\vee,
\]
其中格 $\Lambda^\vee\subset T^*B$ 局部由 $\mathrm{d}x_1,\dots,\mathrm{d}x_n$ 生成，而 $\omega$ 可以与 $T^*B/\Lambda^\vee$ 上的典范辛形式
\[
\omega=\sum_{i=1}^{n}\mathrm{d}x_i\wedge\mathrm{d}u_i
\]
等同。这里 $u_1,\dots,u_n$ 是 $T^*B$ 上的纤维坐标。
% [Tutor Note: 注意 $\Lambda^\vee\subset T^*B$ 与后文 $\Lambda\subset TB$ 的对偶关系：余切丛中的格由 $\mathrm{d}x_i$ 生成，切丛中的格由 $\partial/\partial x_i$ 生成，两者互为对偶格。]

于是由 SYZ 猜想，$X$ 的镜像应当是 $\mu$ 的纤维对偶，即
\[
\check X:=TB/\Lambda,
\]
其中格 $\Lambda\subset TB$ 局部由 $\partial/\partial x_1,\dots,\partial/\partial x_n$ 生成。注意 $\check X$ 自然地是一个复流形，其全纯坐标由 $z_i:=\exp(x_i+\mathbf{i}y_i)$ 给出，其中 $y_1,\dots,y_n$ 是 $TB$ 上与 $u_1,\dots,u_n$ 对偶的纤维坐标。$X$ 上配有全纯体积形式
\[
\check\Omega:=\mathrm{d}\log z_1\wedge\cdots\wedge\mathrm{d}\log z_n .
\]
% [Tutor Note: 原文 p.4 此处印作 "$X$ is equipped with the holomorphic volume form $\check\Omega$"，但按上下文（前一句刚说明 $\check X$ 是复流形、全纯坐标为 $z_i$）应为 $\check X$：500 dpi 裁图 work/tmp/z_p2_4_volform.png 显示同一段上一行的 $\check X$ 带 check，而此处的 $X$ 无 check，确认是原文笔误（漏掉 check）。译文正文照原样保留 $X$，读者应读作 $\check X$。]
% [Tutor Note: 虚数单位 $\mathbf{i}$ 在原文以正体粗体排版（区别于指标 $i$），故 $z_i=\exp(x_i+\mathbf{i}y_i)$ 中的 $\mathbf{i}$ 不应与指标 $i$ 混淆。]

存在一个显式的纤维方向上的 Fourier 型变换，称为半平坦 SYZ 变换 $\mathcal{F}^{\text{semi-flat}}$，它把 $\exp\mathbf{i}\omega$ 送到 $\check\Omega$（例如见 [23, 第 2 节]）。半平坦 SYZ 镜像对称的许多其他细节由 Leung 在 [120] 中完成。

如果我们现在切换到 $B$ 上的复仿射结构，就会在 $\check X$ 上得到一个与其复结构相容的辛结构，使得镜像 $\check X$ 成为 Kähler 流形。此外，若上述函数 $\phi$ 满足实 Monge--Ampère 方程
\[
\det\Bigl(\frac{\partial^2\phi}{\partial x_i\partial x_j}\Bigr)=\text{常数},
\]
则我们在 $X$ 与 $\check X$ 上都得到 $T^n$ 不变的 Ricci 平坦度量。$B$ 上的 McLean 度量此时称为 Monge--Ampère 度量，而 $B$ 称为 Monge--Ampère 流形。这把镜像对称与实 Monge--Ampère 方程及仿射 Kähler 几何的研究联系起来；早在镜像对称被发现之前，Cheng 与 Yau 就已在这一方向作出了重要贡献 [35, 36, 37]。事实证明，在带奇点的仿射流形上构造 Monge--Ampère 度量是一个非常困难的问题。Loftin--Yau--Zaslow 的高度非平凡的工作 [131, 132] 在"Y"型顶点附近构造出了这类度量，而"Y"型顶点是三维情形中一种典型的奇点类型。但除此之外，人们所知甚少。
% [Tutor Note: "real Monge--Ampère equation"（实 Monge--Ampère 方程）指 $\det(\partial^2\phi/\partial x_i\partial x_j)=\text{常数}$，出现在 $B$（实 $n$ 维仿射流形）上；不要与复 Monge--Ampère 方程混淆。]
% [Tutor Note: 原文的 "$\det(\cdot)=\text{constant}$" 是 $\phi$ 的 Hessian 行列式为常数；结合 McLean 度量的定义，这等价于 $T^n$-不变的 Ricci 平坦性。]

% [原文 p.5]
\section{SYZ 纤维化的构造}
在 1996 年 SYZ 猜想提出之后不久，人们投入了大量精力来寻找 Calabi--Yau 流形上特殊拉格朗日环面纤维化（即 SYZ 纤维化）的例子。\footnote{在 [69, 70] 中，Gross 并没有构造这类纤维化，而是假设它们存在，并推导出镜像对称所预言的若干有趣推论。}Zharkov [182] 首先在光滑射影环面簇 $\PP_\Delta$ 中的 Calabi--Yau 超曲面上构造了拓扑环面纤维化，其中包括重要的五次三维流形例子。他通过把 $\PP_\Delta$ 上矩映射在矩多面体边界 $\partial\Delta$ 上的限制形变到附近的光滑 Calabi--Yau 超曲面，得到了这些纤维化。
% [Tutor Note: $\PP_\Delta$ 表示由多面体 $\Delta$ 决定的环面簇（射影）；$\partial\Delta$ 是矩多面体的边界。两者不可混写为 $\PP$ 与 $\Delta$。]

应用类似的想法与梯度-Hamilton 流，W.-D. Ruan 在一系列论文 [151, 152, 153] 中找到了五次三维流形上的拉格朗日环面纤维化。他还对这些纤维化的单值（monodromy）做了非平凡的计算，这些计算后来被 Gross [72] 用来给出 Calabi--Yau 流形上 SYZ 镜像对称的一个拓扑版本。另一方面，Mikhalkin [137] 利用热带几何的工具，在环面簇的超曲面上构造了光滑环面纤维化。

一般而言，构造特殊拉格朗日子流形是非常困难的问题。一个有希望的方法是平均曲率流。Thomas [164] 对 Calabi--Yau 流形中 Maslov 指标为零的拉格朗日子流形类提出了一个稳定性概念，它应当镜像于全纯向量丛的稳定性。Thomas--Yau [165] 猜想：一个 Hamiltonian 同痕类中存在唯一的特殊拉格朗日代表元，当且仅当该类是稳定的；并且这样的代表元应当由平均曲率流得到，而平均曲率流的长时间存在性应当成立。他们进一步对特殊拉格朗日子流形提出了 Jordan--Hölder 型分解，并将其与平均曲率流中奇点的形成联系起来。

这些提法与猜想对 Calabi--Yau 几何与 SYZ 猜想的发展有很大影响，这一领域已取得许多进展：[32, 33, 104, 118, 140, 141, 143, 158, 159, 160, 170, 171, 172, 173]；更多细节见 [174] 与 [142] 及其中的参考文献。遗憾的是，尽管付出了如此多的努力，五次三维流形上特殊拉格朗日环面纤维化的存在性仍然未知。

相比之下，非紧的特殊拉格朗日纤维化的例子要容易得到得多。Harvey 与 Lawson 关于校准几何的著名论文 [94] 给出了其中最简单的例子：由
\[
\begin{aligned}
f:\C^3&\to\R^3,\\
(z_1,z_2,z_3)&\mapsto\bigl(\operatorname{Im}(z_1z_2z_3),\,|z_1|^2-|z_2|^2,\,|z_1|^2-|z_3|^2\bigr),
\end{aligned}
\]
定义的映射是一个特殊拉格朗日纤维化，其纤维在对角 $T^2$-作用之下不变。这一点后来被 Goldstein [65] 与 Gross [71] 的独立工作大大推广，他们在任意环面 Calabi--Yau $n$-重上构造了显式的特殊拉格朗日环面纤维化。这些例子的判别轨迹具有实余维二，并且可以被显式地描述。
% [Tutor Note: 原文 "is a special Lagrangian fibrations"（不定冠词 a 配复数名词）是英文笔误，应为单数 "a special Lagrangian fibration"；译文按单数处理，不改变数学内容。]
% [Tutor Note: 该映射的纤维是 $\C^3$ 在对角 $T^2$-作用下的轨道，$f$ 的分量用到 $\operatorname{Im}(z_1z_2z_3)$（注意原文印刷为 $z_1z_2z_3$ 而非 $z_1z_2\bar z_3$）以及 $|z_1|^2-|z_2|^2$、$|z_1|^2-|z_3|^2$。]
